% =========================================================
\chapter{C05 --- Diagrammes de Bode}
% =========================================================

\begin{rembox}[title=En bref]
\textbf{De quoi parle ce chapitre ?} Dessiner, pour toutes les fréquences, ce que fait une
fonction de transfert : le \textbf{diagramme de Bode en module} (rapport des amplitudes
sortie/entrée) et le \textbf{diagramme de Bode en phase} (déphasage sortie/entrée). \\
\textbf{Les trois objectifs} :
\begin{enumerate}
  \item \textbf{factoriser} les fonctions de transfert en un produit de termes fondamentaux
        du premier ordre (comme on décompose un polynôme de degré N) ;
  \item suivre une \textbf{recette de cuisine systématique} : comportement physique aux
        limites ($\omega\to0$, $\omega\to\infty$), forme canonique à partir des impédances,
        vérification de la cohérence, tracé du module et de l'argument ;
  \item la \textbf{superposition asymptotique} : dessiner chaque terme élémentaire seul (des
        droites), puis \textbf{additionner graphiquement} ces droites, intervalle par
        intervalle.
\end{enumerate}
\end{rembox}

\section{Objectifs : la forme canonique}

\slide{c05-04}{La forme canonique et les règles de dessin.}

\begin{defnbox}[title=Forme canonique (cinq termes du premier ordre)]
\[ \underline H(j\omega) = K\cdot
   \frac{j\frac{\omega}{\omega_{Z1}}\cdot\left(1 + j\frac{\omega}{\omega_{Z2}}\right)\cdot
         \left(1 + j\frac{\omega}{\omega_{Z3}}\right)\cdots\left(1 + j\frac{\omega}{\omega_{Zk}}\right)}
        {j\frac{\omega}{\omega_{P1}}\cdot\left(1 + j\frac{\omega}{\omega_{P2}}\right)\cdot
         \left(1 + j\frac{\omega}{\omega_{P3}}\right)\cdots\left(1 + j\frac{\omega}{\omega_{Pl}}\right)} \]
\begin{itemize}
  \item $K$ est une constante ;
  \item $\omega_{Zi}$ ($i = 1..k$) : les \textbf{zéros} (pulsations caractéristiques au
        \textbf{numérateur}) ;
  \item $\omega_{Pi}$ ($i = 1..l$) : les \textbf{pôles} (au \textbf{dénominateur}).
\end{itemize}
La notation est libre, mais doit permettre de différencier les pulsations. Les cinq termes
de base sont donc : $K$, $j\frac\omega{\omega_Z}$, $\frac1{j\omega/\omega_P}$,
$1 + j\frac\omega{\omega_Z}$, $\frac1{1 + j\omega/\omega_P}$.
\end{defnbox}

\begin{defnbox}[title=Règles de dessin]
\begin{itemize}
  \item Le \textbf{module} $|\underline H(j\omega)|$ est exprimé en \textbf{dB} (décibels),
        sur une échelle \textbf{linéaire} ;
  \item l'\textbf{argument} $\mathrm{Arg}(\underline H(j\omega))$ en \textbf{radians ou
        degrés}, échelle \textbf{linéaire} ;
  \item dans les deux cas, l'axe des pulsations est sur une \textbf{échelle
        logarithmique}.
\end{itemize}
\end{defnbox}

\section{Rappel : différence entre R et Z}

\slide{c05-05}{Diviseur résistif : le rapport est une valeur. Diviseur d'impédances : une fonction de $\omega$.}

\begin{propbox}[title={Diviseur résistif vs fonction de transfert}]
\begin{itemize}
  \item Avec deux résistances : $\frac{u_S}{u} = \frac{R_2}{R_1 + R_2}$ : \textbf{le rapport
        est une valeur}. Quelle que soit la fréquence, l'amplitude de sortie est un rapport
        constant de l'entrée, et il n'y a pas de décalage.
  \item Avec des impédances : $\frac{\underline u_S}{\underline u} = \frac{\underline Z_2}{\underline Z_1 + \underline Z_2} = \underline H(j\omega)$ :
        \textbf{le rapport est une fonction}, la fonction de transfert, avec un
        \textbf{module} (l'amplitude varie avec la fréquence) et un \textbf{argument} (la
        sortie est décalée).
\end{itemize}
\end{propbox}

\section{Pourquoi des produits, et pourquoi des dB ?}

\slide{c05-06}{Les modules en dB des termes de base, et l'origine du facteur 20.}

